\chapter{Git 与可复现科研}

\section{本章阅读方式}

本章快读 Git 概念，重点学会把一次科研结果冻结到可定位版本。评分重点不是 commit 数量，而是别人能否知道哪个版本生成了哪个图、表或 record。

最低交付是一个 Evidence Record，以及 \texttt{git log}、\texttt{git status} 或 tag 记录。本章进入 Part III 的作用，是让机器学习实验、数据处理结果和报告图表都能追踪到具体代码版本。

\section{本章导读：Git 不是备份工具，而是科研时间轴}

很多学生第一次接触 Git 时，会把它理解成“保存代码版本的工具”。这个理解没有错，但太窄。对于科研项目来说，Git 更重要的作用是建立一条可追踪的时间轴：什么时候改了数据筛选条件，什么时候修正了图表脚本，什么时候生成了报告中的关键结果，哪个 commit 对应期末提交版本。

如果没有版本控制，科研项目很容易变成一堆文件名：

\begin{lstlisting}[style=terminal]
analysis.py
analysis_new.py
analysis_final.py
analysis_final_really.py
figure3_old.png
figure3_fixed.png
report_latest.docx
\end{lstlisting}

这些文件名看似记录了历史，实际上几乎无法回答关键问题：为什么改，改了什么，哪一个版本生成了最终结果，错误从哪一步开始出现。Git 的价值就在于把“文件变化”变成可以检查、可以讨论、可以回到过去的项目历史。

\begin{defbox}[版本控制]
版本控制是一种系统记录文件变化的方法。它不仅保存历史，还帮助研究者回答“谁在什么时候、因为什么理由、对哪些文件做了什么修改”。
\end{defbox}

本章会把 Git 当作科研记录系统来讲，而不是命令清单。你会学习命令，但每条命令都要回到一个科研问题：这一步如何帮助结果可复查、协作可管理、错误可定位、发布可冻结。

\section{学习目标与训练目标}

学完本章后，你应该能够：

\begin{itemize}
    \item 解释为什么版本控制是科研工作的一部分，而不是程序员专用技能；
    \item 理解 working directory、staging area 和 repository 的分工；
    \item 使用 \texttt{git status} 和 \texttt{git diff} 判断当前项目状态；
    \item 使用 \texttt{git add} 和 \texttt{git commit} 创建主题清楚的提交；
    \item 使用 \texttt{git log}、branch、merge、tag 和 remote 理解项目历史；
    \item 知道如何处理最基本的 merge conflict；
    \item 区分应该进 Git 的文本记录和不应该直接进 Git 的大型数据、自动生成结果；
    \item 把最终结果、运行环境、数据来源和 commit hash 连接成最小 evidence record。
\end{itemize}

训练材料的目标是让你在临时目录中真正操作一个小型仓库：初始化、修改文件、查看 diff、提交、创建分支、模拟冲突、打 tag，并生成一份版本化结果记录。

\section{一个贯穿本章的小型科研情景}

假设你正在做一个很小的光度测量练习。项目目录中有：

\begin{lstlisting}[style=terminal]
mini_photometry/
  README.md
  scripts/
    measure_flux.py
  results/
    flux_summary.csv
  notes/
    interpretation.md
\end{lstlisting}

一开始，脚本只是粗略计算目标亮度。后来你发现背景扣除公式写错了，又修改了结果表和解释说明。再后来，同学也修改了 \texttt{README.md}，你们需要合并。最后，期末报告要引用一个冻结版本。

这个情景包含 Git 的所有核心问题：

\begin{itemize}
    \item 哪些修改还没有进入历史？
    \item 这次提交到底应该包含哪些文件？
    \item 背景扣除修正前后改了什么？
    \item 哪个版本生成了最终 \texttt{flux\_summary.csv}？
    \item 大型原始图像是否应该直接提交？
    \item 如何把期末报告对应的版本冻结下来？
\end{itemize}

接下来所有命令都围绕这些问题展开。

\section{Git 的三层模型}

Git 最容易让初学者困惑的地方，是文件修改不会自动进入历史。Git 把项目状态分成三个层次，见表~\ref{tab:ch02_git_layers}：

\begin{table}[htbp]
\centering
\begin{tabular}{p{0.25\textwidth}p{0.28\textwidth}p{0.35\textwidth}}
\toprule
层次 & 保存什么 & 科研含义 \\
\midrule
Working directory & 当前正在编辑的文件 & 实验台面，可能混有多个未完成想法 \\
Staging area & 准备进入下一次提交的修改 & 本次要记录的主题，需要主动选择 \\
Repository & 已提交的历史快照 & 项目时间轴，可复查、可回到过去 \\
\bottomrule
\end{tabular}
\caption{Git 工作流中的三层项目状态及其科研含义。}
\label{tab:ch02_git_layers}
\end{table}

\begin{examplebox}[最小 Git 工作流]
\begin{lstlisting}[style=terminal]
edit files
git status
git diff
git add selected_files
git diff --staged
git commit -m "Explain the change"
\end{lstlisting}
\end{examplebox}

staging area 的存在不是麻烦，而是帮助你把混乱的工作区整理成清楚的历史节点。科研中你常常同时做了多件事：改图表标题、修 bug、尝试新参数、写报告说明。它们不一定应该进入同一个 commit。一次好的提交应该尽量对应一个明确理由。

\section{Commit、HEAD、branch 与 tag 的心智模型}

在开始记命令之前，先把 Git 的几个核心名字放到一张脑内地图里，会少走很多弯路。表~\ref{tab:ch02_commit_names} 先给出这些名字在科研项目中的角色。Git 仓库里的历史不是一串文件名，而是一串 commit。每个 commit 可以粗略理解为：

\begin{center}
一次项目快照 + 作者和时间 + 提交说明 + 指向上一个 commit 的线索。
\end{center}

因此，commit 不是“只保存改动的一句话”，而是一个可以被重新找到的历史节点。你在报告里写“本图由 commit \texttt{3a7c91e} 生成”，意思就是读者可以回到那一刻的项目状态，检查脚本、参数和说明是否一致。

\begin{table}[htbp]
\centering
\begin{tabular}{p{0.20\textwidth}p{0.30\textwidth}p{0.36\textwidth}}
\toprule
名字 & 可以怎样理解 & 科研工作中的作用 \\
\midrule
commit & 一个带说明的项目历史节点 & 说明一次结果或一次修正对应的代码状态 \\
HEAD & 当前所在的历史位置 & 告诉你现在正在基于哪一版继续工作 \\
branch & 会随着新提交移动的名字 & 让稳定主线和试验想法分开推进 \\
tag & 固定贴在某个 commit 上的名字 & 冻结报告、作业或 release 对应版本 \\
\bottomrule
\end{tabular}
\caption{Commit、HEAD、branch 与 tag 的最小心智模型。}
\label{tab:ch02_commit_names}
\end{table}

一个常见误解是把 branch 和 tag 都理解成“版本”。它们确实都指向 commit，但行为不同：branch 会随着你继续提交而向前移动，tag 通常固定不动。主线 \texttt{main} 像一支还在继续写的实验记录本；\texttt{v0.1-analysis} 这样的 tag 像给某一页贴了书签，说明这一页可以被引用。

这也解释了为什么最终结果不能只写“来自 main 分支”。\texttt{main} 明天可能已经前进了。更稳妥的记录是 commit hash 或 release tag。它们把结果固定到一个明确历史节点，而不是一个会移动的名字。

\section{读懂 \texttt{git status} 后该做什么}

\texttt{git status} 不只是告诉你“有文件改了”，还在提示下一步。可以把常见状态翻译成表~\ref{tab:ch02_status_decisions} 这样的决策表：

\begin{table}[htbp]
\centering
\begin{tabular}{p{0.28\textwidth}p{0.34\textwidth}p{0.26\textwidth}}
\toprule
状态提示 & 含义 & 推荐动作 \\
\midrule
\texttt{modified} & 已跟踪文件被修改，还未暂存 & 先看 \texttt{git diff} \\
\texttt{untracked} & Git 还不知道这个新文件 & 判断是否应进入 Git \\
\texttt{changes to be committed} & 已暂存，准备提交 & 看 \texttt{git diff --staged} \\
\texttt{nothing to commit} & 工作区和暂存区干净 & 可以运行、拉取或切分支 \\
\texttt{your branch is ahead} & 本地有远程没有的提交 & 检查后可以 push \\
\bottomrule
\end{tabular}
\caption{\texttt{git status} 常见状态提示与推荐下一步动作。}
\label{tab:ch02_status_decisions}
\end{table}

这张表训练的是 Git 的“状态机”意识。不要把 Git 命令当成咒语；每次先读状态，再决定是查看差异、暂存、提交、忽略，还是推送。

\section{初始化仓库与第一次提交}

在项目根目录中运行：

\begin{examplebox}[初始化仓库]
\begin{lstlisting}[style=terminal]
$ git init
$ git status
\end{lstlisting}
\end{examplebox}

\texttt{git init} 会创建隐藏目录 \texttt{.git/}，里面保存仓库元数据和历史。不要手动编辑 \texttt{.git/}，把它当作 Git 自己维护的数据库即可。

第一次提交通常应该包含项目的最小说明、脚本入口和忽略规则：

\begin{examplebox}[第一次提交]
\begin{lstlisting}[style=terminal]
$ git add README.md scripts/measure_flux.py .gitignore
$ git status
$ git commit -m "Create minimal photometry project"
\end{lstlisting}
\end{examplebox}

注意这里没有使用 \texttt{git add .}。入门阶段更推荐显式列出文件，因为它迫使你确认本次提交到底包含什么。

\section{\texttt{git status}：先看状态，不要靠猜}

\texttt{git status} 是 Git 中最重要的诊断命令。它回答：

\begin{itemize}
    \item 当前在哪个分支；
    \item 哪些文件被修改；
    \item 哪些修改已经 staged；
    \item 哪些文件还没有被跟踪；
    \item 是否与远程仓库有同步差异。
\end{itemize}

\begin{protip}[在 Git 里迷路时先看 status]
当你不确定当前状态时，不要猜，不要随便 pull、merge 或 restore。先运行 \texttt{git status}，读懂它再决定下一步。
\end{protip}

科研项目中，\texttt{status} 的意义类似实验记录中的“当前状态检查”。它不改变文件，只帮助你确认下一步是否安全。

\section{\texttt{git diff}：提交前的审稿}

\texttt{git diff} 显示工作区和上一次提交之间的差异。它回答的问题是：如果我现在提交，这些变化到底是什么？

\begin{examplebox}[查看差异]
\begin{lstlisting}[style=terminal]
$ git diff
$ git diff scripts/measure_flux.py
\end{lstlisting}
\end{examplebox}

科研中阅读 diff 时，要特别关注：

\begin{itemize}
    \item 是否改了数据筛选条件；
    \item 是否改了随机种子、阈值、单位转换或背景扣除公式；
    \item 是否混入临时调试输出；
    \item 是否意外修改了报告、图表或配置文件；
    \item 是否包含绝对路径、私有路径或本地机器特有设置。
\end{itemize}

已经 staged 的内容可以用：

\begin{examplebox}[查看将要提交的差异]
\begin{lstlisting}[style=terminal]
$ git diff --staged
\end{lstlisting}
\end{examplebox}

\texttt{git diff --staged} 相当于提交前的最后审稿。它帮助你确认下一次 commit 是否主题清楚、边界明确。

\section{提交：把一次修改变成可解释的历史节点}

\texttt{git add} 选择要进入下一次提交的修改，\texttt{git commit} 创建历史快照。

\begin{examplebox}[一次主题清楚的提交]
\begin{lstlisting}[style=terminal]
$ git add scripts/measure_flux.py notes/interpretation.md
$ git diff --staged
$ git commit -m "Fix background subtraction in flux estimate"
\end{lstlisting}
\end{examplebox}

一个科研 commit message 不应该只是 \texttt{update}、\texttt{fix} 或 \texttt{final}。表~\ref{tab:ch02_commit_messages} 给出了一组从弱信息到可复查信息的对照。它最好说明：

\begin{itemize}
    \item 改了什么；
    \item 为什么改；
    \item 这次改动对应哪个问题、检查或结果解释。
\end{itemize}

\begin{table}[htbp]
\centering
\begin{tabular}{p{0.38\textwidth}p{0.44\textwidth}}
\toprule
较弱的提交信息 & 更好的提交信息 \\
\midrule
\texttt{update} & \texttt{Add input file count check} \\
\texttt{fix bug} & \texttt{Fix background subtraction sign} \\
\texttt{plot} & \texttt{Label flux residual plot axes} \\
\texttt{final version} & \texttt{Freeze report figures for draft submission} \\
\bottomrule
\end{tabular}
\caption{较弱提交信息与更可复查的科研提交信息对照。}
\label{tab:ch02_commit_messages}
\end{table}

好的 commit message 是给未来的自己、同学和助教看的。几周后你不一定记得当时为什么改，但项目历史应该能提醒你。

\section{一次提交应该有边界}

假设你同时做了三件事：

\begin{itemize}
    \item 修正 \texttt{measure\_flux.py} 中背景扣除符号；
    \item 改写 README 中的项目介绍；
    \item 临时试验一个新的孔径半径，生成了测试图。
\end{itemize}

这三件事最好不要机械塞进同一个提交。更合理的做法是：

\begin{enumerate}
    \item 先提交背景扣除 bug fix，因为它改变结果含义；
    \item 再提交 README 说明，因为它是文档更新；
    \item 临时测试图若不能复现或只是探索输出，可以不提交，或先放进被忽略目录。
\end{enumerate}

提交边界的原则是：\textbf{一次提交应该能用一句话说明为什么存在。} 如果一句话里出现了“顺便”“还有”“以及很多东西”，通常说明这次提交太杂。

\section{撤销与恢复：先分清要撤销哪一层}

Git 中“撤销”不是一个动作，而是一组不同层次的操作。入门阶段先掌握两个安全边界：

\begin{examplebox}[取消暂存但保留文件内容]
\begin{lstlisting}[style=terminal]
$ git restore --staged README.md
\end{lstlisting}
\end{examplebox}

这条命令把 \texttt{README.md} 从 staging area 拿出来，但工作区中的实际修改还保留。它适合处理“这个文件不属于本次提交”的情况。

\begin{examplebox}[丢弃工作区修改]
\begin{lstlisting}[style=terminal]
$ git restore README.md
\end{lstlisting}
\end{examplebox}

这条命令会把 \texttt{README.md} 恢复到上一次提交的状态，未提交修改会消失。使用前一定要先运行 \texttt{git diff}。

\begin{warnbox}[恢复命令不是回收站]
\texttt{git restore} 可以帮你撤销错误修改，也可能直接丢掉尚未提交的工作。只要不确定，先看 \texttt{git status} 和 \texttt{git diff}，必要时把当前文件复制到临时位置再恢复。
\end{warnbox}

\section{查看历史：从文件列表回到项目故事}

\texttt{git log} 显示提交历史。常用形式是：

\begin{examplebox}[查看紧凑历史]
\begin{lstlisting}[style=terminal]
$ git log --oneline --graph --decorate
$ git show --stat HEAD
\end{lstlisting}
\end{examplebox}

\texttt{git log} 让你看到项目如何变化，\texttt{git show --stat HEAD} 则显示最近一次提交改了哪些文件。科研中它们可以帮助你回答：

\begin{itemize}
    \item 某个异常结果是从哪次修改后出现的；
    \item 某个图表脚本是否已经修正；
    \item 报告中的结果对应哪个提交；
    \item 期末提交版本是否已经被冻结。
\end{itemize}

如果提交信息写得清楚，\texttt{git log} 就像项目日记；如果提交信息全是 \texttt{update}，历史就会变成噪声。

\section{分支：把试验和稳定主线分开}

科研工作经常需要试错。例如你想尝试新的背景估计、新的特征工程、新的模型或新的图表风格。直接在主线修改当然可以，但如果实验失败，主线可能被打乱。分支让你在不破坏稳定版本的前提下试验。

\begin{examplebox}[创建试验分支]
\begin{lstlisting}[style=terminal]
$ git switch -c experiment/robust-background
$ git status
\end{lstlisting}
\end{examplebox}

试验完成后，可以回到主线并合并：

\begin{examplebox}[合并分支]
\begin{lstlisting}[style=terminal]
$ git switch main
$ git merge experiment/robust-background
\end{lstlisting}
\end{examplebox}

分支不是为了复杂化项目，而是为了让“可工作的主线”和“正在试的想法”分开。对于本科项目，一个简单原则就够了：主线尽量保持能运行，实验性修改放到分支。

\section{冲突：Git 需要人做语义判断}

merge conflict 的本质是：Git 发现两个版本修改了同一个位置，它无法自动判断哪个版本才符合人的意图。这不是 Git 坏了，而是项目出现了需要人理解内容的地方。

冲突文件中会出现由小于号、等号和大于号组成的边界标记。为了避免把教材示例误识别成真实冲突残留，这里用文字替代那些符号：

\begin{lstlisting}[style=terminal]
[current branch version begins]
current version
[incoming version begins]
incoming version
[incoming version ends]
\end{lstlisting}

解决冲突的基本流程是：

\begin{enumerate}
    \item 运行 \texttt{git status} 找到冲突文件；
    \item 打开文件，理解两边改动；
    \item 手动编辑成正确内容，删除冲突标记；
    \item 运行必要检查或测试；
    \item \texttt{git add} 冲突文件；
    \item 完成 merge commit。
\end{enumerate}

\begin{warnbox}[不要机械选择一边]
冲突解决不是简单保留上半段或下半段。尤其是报告、参数、数据筛选和解释文字发生冲突时，需要根据科研含义判断最终内容。
\end{warnbox}

\section{远程仓库：协作和备份不是同一件事}

本地 Git 记录你自己的历史，远程仓库让多人共享同一条历史。常见平台包括 GitHub、GitLab 和学校内部 Git 服务。

\begin{examplebox}[远程仓库基本命令]
\begin{lstlisting}[style=terminal]
$ git clone <repo_url>
$ git pull origin main
$ git push origin main
\end{lstlisting}
\end{examplebox}

\texttt{clone} 是下载一份仓库，\texttt{pull} 是拉取远程更新并整合到本地，\texttt{push} 是上传本地提交。协作项目中，push 前最好先确认：

\begin{itemize}
    \item 本地工作区是干净的；
    \item 最近提交主题清楚；
    \item 不包含私有 token、大型数据或无关临时文件；
    \item 已经 pull 过远程最新变化；
    \item 必要检查能通过。
\end{itemize}

远程仓库也可以作为备份，但它首先是协作历史。不要把它当作网盘随便塞文件。

\section{Tag 与 release：冻结可引用版本}

科研报告、课程作业和公开发布都需要一个稳定版本。只说“用最新代码生成”是不够的，因为最新代码会继续变化。tag 可以给某个提交贴上稳定标签。

\begin{examplebox}[创建分析版本标签]
\begin{lstlisting}[style=terminal]
$ git tag -a v0.1-analysis -m "Freeze first reproducible analysis"
$ git log --oneline --decorate -n 3
\end{lstlisting}
\end{examplebox}

release 是比 tag 更完整的发布单元。一个课程项目 release 至少应该包含：

\begin{itemize}
    \item Git tag 或 commit hash；
    \item 环境说明；
    \item 数据来源和版本说明；
    \item 生成关键结果的脚本或 notebook 入口；
    \item 最终图表和报告；
    \item 已知限制和未解决问题。
\end{itemize}

这和本书后面的 项目证据卡 是同一思想：结果不是孤立文件，而是一组可复查证据。

\section{什么应该进 Git，什么不应该}

Git 擅长管理文本文件，例如源代码、脚本、LaTeX、Markdown、配置文件、小型 CSV 示例和 README。它不适合直接管理大型原始数据、频繁变化的二进制输出和自动生成的缓存。表~\ref{tab:ch02_git_include_exclude} 给出最小判断边界。

\begin{table}[htbp]
\centering
\begin{tabular}{p{0.30\textwidth}p{0.48\textwidth}}
\toprule
通常应该进入 Git & 通常不直接进入 Git \\
\midrule
脚本、notebook 源文件、README、环境文件、小型教学数据、数据清单、报告源文件 & 大型 FITS 原始数据、大型模型权重、自动生成图像、临时缓存、个人密钥、下载目录 \\
\bottomrule
\end{tabular}
\caption{科研项目中通常应该进入 Git 与通常不直接进入 Git 的文件类型。}
\label{tab:ch02_git_include_exclude}
\end{table}

\texttt{.gitignore} 用来告诉 Git 忽略哪些文件：

\begin{examplebox}[一个最小科研项目 \texttt{.gitignore}]
\begin{lstlisting}[style=terminal]
__pycache__/
.ipynb_checkpoints/
data/raw/
results/*.png
results/*.fits
logs/*.tmp
.env
\end{lstlisting}
\end{examplebox}

数据不进 Git，不等于数据不可追踪。大型数据本体可以放在仓库外，但数据来源、下载方式、校验信息、字段说明和生成脚本应该进入 Git。

\begin{examplebox}[数据说明应该进入 Git]
\begin{lstlisting}[style=terminal]
data/
  README.md
  manifest.csv
  checksums.txt
scripts/
  download_demo_data.py
\end{lstlisting}
\end{examplebox}

这条边界非常重要：Git 不一定保存所有数据本体，但应该保存足够信息，让别人知道数据从哪里来、如何验证、如何重新生成。

\section{大型数据如何留下版本线索}

如果原始 FITS 文件太大，不适合进入 Git，仍然可以记录它的身份。最常见的低成本方法是保存文件大小、来源和校验和：

\begin{examplebox}[生成校验线索]
\begin{lstlisting}[style=terminal]
$ shasum -a 256 data/raw/night1_target042.fits \
    > data/checksums.txt
$ ls -lh data/raw/night1_target042.fits \
    >> data/manifest_snapshot.txt
\end{lstlisting}
\end{examplebox}

校验和不是科学解释，但它能帮助你确认“这份文件是否还是同一份文件”。如果同名文件被重新下载、裁剪或覆盖，校验和通常会改变。对于课程项目，哪怕不使用完整数据版本系统，保存 manifest 和 checksum 也能显著提高复查能力。

\section{Notebook 与 Git：友好界面和清晰 diff 的矛盾}

Jupyter notebook 对探索和教学很友好，但它和 Git 的配合并不完美。notebook 文件本质上是 JSON，里面可能包含代码、Markdown、输出、图像和执行计数。一个很小的输出变化，也可能让 diff 变得很长。

在科研项目中，notebook 使用至少应遵守三条底线：

\begin{itemize}
    \item 从头到尾重新运行应该能得到一致结果；
    \item 关键参数、随机种子和输入路径要显式写出；
    \item 最终结果最好由脚本或清晰 notebook 入口生成，而不是依赖隐藏的交互状态。
\end{itemize}

如果一个 notebook 只能在你当前内存状态下跑通，或者必须依赖乱序执行，那么 Git 只能记录文件变化，不能保证科学可复现。

\section{把结果和代码版本连起来}

Git 对科研最有价值的地方，常常不是“可以回退”，而是能把一次结果和一个明确代码版本连起来。一个图表、模型分数或报告结论，如果不知道来自哪个 commit，就很难复查。

\begin{examplebox}[记录当前版本]
\begin{lstlisting}[style=terminal]
$ git rev-parse --short HEAD
$ git log --oneline -5
\end{lstlisting}
\end{examplebox}

对于最终图表和最终报告，至少应该记录：

\begin{itemize}
    \item 生成结果的 commit hash；
    \item 使用的数据版本或数据清单；
    \item 运行入口，例如脚本或 notebook；
    \item 关键参数和随机种子；
    \item 输出文件路径；
    \item 已知限制。
\end{itemize}

\begin{examplebox}[最小版本化结果记录]
\begin{lstlisting}[style=terminal]
Result: results/flux_summary.csv
Commit: 3a7c91e
Data source: data/README.md, manifest row demo-night1
Command: python scripts/measure_flux.py
Environment: environment.yml
Key parameters: aperture_radius=5, background_annulus=8:12
Checks: row count = 20, no negative flux after background subtraction
Limit: simulated teaching data, not calibrated science product
\end{lstlisting}
\end{examplebox}

这份记录使用统一 Evidence Record 母模板，本章只补充以下版本字段：

\begin{examplebox}[Ch02 Evidence Record 附加字段]
\begin{lstlisting}[style=terminal]
Ch02 Fields:
- repository root:
- relevant commit:
- tag/checkpoint:
- tracked files:
- ignored files:
- output linked to commit:
\end{lstlisting}
\end{examplebox}

本章的特殊通过条件是：必须能指出哪个 commit 或 tag 生成了哪个输出文件。这份记录会在后续章节继续升级，最终成为 Evidence Record 和 Model Experiment Record 的一部分。

\section{从日常提交到 release 的算法}

Git 日常提交解决的是“每一步变化可追踪”，release 解决的是“某个阶段结果可引用”。把两者连起来，可以形成一条固定算法：

\begin{enumerate}
    \item 运行项目检查，例如脚本、notebook 或 LaTeX 编译；
    \item 生成或更新关键结果；
    \item 写 evidence record，记录数据、命令、环境和限制；
    \item 用 \texttt{git status} 确认只包含应提交内容；
    \item 用 \texttt{git diff} 和 \texttt{git diff --staged} 审查；
    \item 创建主题清楚的 commit；
    \item 对阶段性版本打 tag；
    \item 在 release note 中说明这个版本能支持什么、不能支持什么。
\end{enumerate}

这条算法会在后面的书稿发布、notebook smoke test、数据 manifest 和 project package 中反复出现。Git 的最终目标不是让历史变复杂，而是让阶段性结果可以被可靠引用。

\section{配套 notebook 如何训练}

本章配套 notebook 会在临时目录中创建一个小型 Git 仓库，不会污染当前项目。训练路线是：

\begin{enumerate}
    \item 初始化仓库；
    \item 创建 README、脚本和结果文件；
    \item 使用 \texttt{status} 和 \texttt{diff} 观察修改；
    \item 做第一次提交；
    \item 修改脚本并提交修正；
    \item 创建分支模拟试验；
    \item 记录当前 commit hash；
    \item 生成最小版本化结果记录。
\end{enumerate}

正文负责讲清楚 Git 为什么是科研时间轴；notebook 负责让你亲手经历“修改、检查、提交、复查”的循环。

\section{常见错误与诊断}

\begin{warnbox}[Git 入门中的典型问题]
\begin{itemize}
    \item 不看 \texttt{git status}，不知道当前处于什么状态；
    \item 不看 \texttt{git diff}，直接 \texttt{git add .}；
    \item commit message 全是 \texttt{update}、\texttt{fix}、\texttt{final}；
    \item 把大型原始数据、自动生成结果或私有密钥提交到仓库；
    \item notebook 乱序执行，结果依赖隐藏状态；
    \item 冲突时机械保留某一边，不理解科研含义；
    \item 最终结果没有记录 commit hash、数据版本和运行入口。
\end{itemize}
\end{warnbox}

诊断 Git 问题时，优先运行三条命令：

\begin{lstlisting}[style=terminal]
$ git status
$ git diff
$ git log --oneline --decorate -n 5
\end{lstlisting}

它们分别回答当前状态、当前变化和最近历史。大多数入门问题都能从这三条命令开始定位。

\section{AI 助手如何帮助你}

AI 助手可以在几个具体环节提供帮助：

\begin{itemize}
    \item 解释 \texttt{git status} 输出；
    \item 总结 \texttt{git diff} 中真正改变的内容；
    \item 建议更清楚的 commit message；
    \item 设计适合科研项目的 \texttt{.gitignore}；
    \item 整理 release checklist；
    \item 帮助阅读冲突文件。
\end{itemize}

但它不能替你判断科研语义。尤其是参数、筛选条件、解释文字和结果文件发生变化时，最终判断必须回到项目问题本身。

比较好的提问方式是：

\begin{lstlisting}[style=terminal]
请阅读下面的 git diff，帮我判断这次修改是否只包含一个主题。
请指出是否混入临时路径、调试输出、数据筛选条件变化或
不应该提交的大文件。
\end{lstlisting}

或者：

\begin{lstlisting}[style=terminal]
这是我的项目目录。请建议一份适合科研项目的 .gitignore，
并解释哪些文件应该进入 Git，哪些文件只需要写入 data README。
\end{lstlisting}

注意，AI 给出的 Git 命令也需要你理解后再运行。尤其是 \texttt{reset}、\texttt{restore}、\texttt{clean}、force push 这类命令，可能丢失工作或改写协作历史，不能盲目复制。

\section{练习与交付}

\begin{exercisebox}[基础练习]
在临时目录中创建一个小型仓库，包含 \texttt{README.md}、\texttt{scripts/measure\_flux.py} 和 \texttt{notes/interpretation.md}。运行 \texttt{git init}、\texttt{git status}，完成第一次提交。
\end{exercisebox}

\begin{exercisebox}[过程练习]
修改脚本中的一个参数或说明文字。依次运行 \texttt{git status}、\texttt{git diff}、\texttt{git add}、\texttt{git diff --staged}、\texttt{git commit}。解释每一步输出回答了什么问题。
\end{exercisebox}

\begin{exercisebox}[诊断练习]
设计一份 \texttt{.gitignore}，让 \texttt{data/raw/}、\texttt{results/*.png}、\texttt{\_\_pycache\_\_/} 不进入 Git，同时保留 \texttt{data/README.md} 和 \texttt{scripts/}。解释为什么数据不进 Git 仍然需要数据说明进 Git。
\end{exercisebox}

\begin{exercisebox}[交付练习]
生成一份最小版本化结果记录，至少包含结果文件、commit hash、数据来源、运行入口、关键参数、检查结果和已知限制。把这份记录保存为 \texttt{notes/evidence\_record.md} 并提交。
\end{exercisebox}

\section{Exit Gate}

\begin{enumerate}
    \item 仓库当前状态是否干净？
    \item 关键代码和 record 是否已 commit？
    \item 是否有可评分 tag 或 commit？
    \item 输出文件能否追踪到版本？
\end{enumerate}

最低交付：Evidence Record 和 \texttt{git log/status/tag} 记录。

\section{本章小结}

Git 的核心不是命令本身，而是把科研项目变成可追踪时间轴。一个可靠项目不只是“文件还在”，而是能说明每个关键结果来自哪一版代码、哪份数据、哪个运行入口和哪些参数。只要你养成 \texttt{status -> diff -> add -> diff --staged -> commit} 的循环，并用 \texttt{.gitignore}、tag、release 和 evidence record 管住数据边界与结果版本，Git 就会从额外负担变成科研可信度的基础设施。
